\documentclass[spanish, 12pt]{article}
\usepackage[T1]{fontenc} % https://tex.stackexchange.com/questions/664/why-should-i-use-usepackaget1fontenc
\usepackage[spanish,es-tabla]{babel} % Traducciones, y cambia cuadro por tabla
\usepackage{csquotes} % Mejor soporte de comillas en bibliografía
\usepackage{circuitikz} % Para armar circuitos directamente en LaTeX
\usepackage{amsmath} % Mejoras en fórmulas matemáticas
\usepackage{multirow} % Multples filas 
\usepackage{multicol} % Múltiples columnas
\usepackage[a4paper, portrait, margin=2.5cm]{geometry} % Ajusto márgenes y tamaño de página
\usepackage{booktabs} % Mejoras en tablas
\usepackage{graphicx} % Extended interface for graphics inclusion
\usepackage{float} % Mejor soporte de float para figuras
\usepackage{cancel} % Permite cancelar términos
\usepackage{ulem} % Subrayados
\usepackage{pdflscape} % Páginas horizontales
\usepackage{mfirstuc} % Primera letra mayúscula

% --- Para listar código
\usepackage{listings} 
\usepackage{xcolor}
% --- Para figuras y subfiguras
\usepackage{caption}
\usepackage{subcaption}
% --- Para el título
\usepackage[absolute,overlay]{textpos}
% --- Para hipervínculos
\usepackage{hyperref}
\hypersetup{
    hidelinks = true
}

\urlstyle{same}

% --- Colores para listados de código
\definecolor{codegreen}{rgb}{0,0.6,0}
\definecolor{codegray}{rgb}{0.5,0.5,0.5}
\definecolor{codepurple}{rgb}{0.58,0,0.82}
\definecolor{backcolour}{rgb}{0.95,0.95,0.92}

\lstdefinestyle{mystyle}{
    commentstyle=\color{codegreen},
    keywordstyle=\color{magenta},
    numberstyle=\tiny\color{codegray},
    stringstyle=\color{codepurple},
    basicstyle=\ttfamily\footnotesize,
    breakatwhitespace=false,
    breaklines=true,
    captionpos=b,
    keepspaces=true,
    numbers=left,
    numbersep=5pt,
    showspaces=false,
    showstringspaces=false,
    showtabs=false,
    tabsize=2
}

\lstset{style=mystyle}
% ---

% --- Configuración de bibiliografía
\usepackage[
backend=biber,
style=alphabetic,
sorting=ynt
]{biblatex}
\addbibresource{references.bib}
% ---

% --- Para informes de tipo pregunta/respuesta
\newcounter{question}
\setcounter{question}{0}
\newcommand{\question}[1]{
    \stepcounter{question}
    \vspace{0.2cm}
    \noindent\rule{\textwidth}{0.4mm}
    {\large\thequestion. #1\par}
    \noindent\rule{\textwidth}{0.4mm}
}
\newcommand{\answer}[1]{
    \vspace{0.2cm}
    #1
}

% --- Unidades
\usepackage{siunitx}
\sisetup{
    separate-uncertainty=true,
    locale = DE
}

% --- Para TODOs y comentarios
\newcommand{\TODO}[1]{
    \begin{center}
        \colorbox{orange}{{\Huge TODO: }#1}
    \end{center}
}

\newcommand{\comment}[1]{
    \begin{center}
        \colorbox{yellow}{#1}
    \end{center}
}

\begin{document}

    \newcommand{\materia}[1]{{\huge #1}\\\vspace{0.5cm}}
\newcommand{\titulo}[1]{{\Huge \textbf{#1}}\\\vspace{0.5cm}}
\newcommand{\subtitulo}[1]{{\huge #1}\\\vspace{0.5cm}}
\newcommand{\autor}[3]{#1\\{\small Mat. #2\\\href{mailto:#3}{#3}}\\\vspace{0.7cm}}
\newcommand{\fecha}[4]{{\small Fecha de #1: #4-#3-#2}\\} % YYYY-MM-DD
\newcommand{\departamento}[1]{{\small Departamento de #1\\Facultad de Ingeniería\\Universidad Nacional de Mar del Plata}}
\newenvironment{mytitlepage}
    {\pagenumbering{gobble}\begin{center}}
    {\end{center}\newpage\pagenumbering{arabic}}

\begin{mytitlepage}
    \includegraphics[width=4cm]{img/logo-fi.pdf}
   
    \begin{textblock*}{\paperwidth}(0cm, 7cm)
    \materia{Física Experimental}
    \titulo{Trabajo Práctico 6}
    \subtitulo{Péndulo de varilla 2}
    \end{textblock*}
    
    \begin{textblock*}{\paperwidth}(0cm, 12cm)
    \autor{Nazareno Montesoro}{13.668}{n.montesoro@gmail.com}
    \autor{Sofía Romero Ortueta}{13.778}{ssofi.romero@gmail.com}
    \end{textblock*}
    
    \begin{textblock*}{\paperwidth}(0cm, 25cm) 
    \fecha{realización}{28}{10}{2022}
    \fecha{entrega}{03}{11}{2022}
    \end{textblock*}
    
    \begin{textblock*}{\paperwidth}(0cm, 27cm) 
    \departamento{Física}
    \end{textblock*}
\end{mytitlepage}

    \tableofcontents

    \section{Introducción y fundamentos}

En este práctico Ud. determinará el valor de la constante de la gravedad
utilizando un péndulo físico de varilla rígida. El período de la oscilación $T$
está relacionado con la longitud $L$ de la varilla y la constante de la
gravedad $g$ por medio de la siguiente ecuación:

\begin{equation}
    \label{ec:intro:periodo}
    T = 2\pi \sqrt{ \frac{2L}{3g} }
\end{equation}

La medida del período $T$ estará afectada por errores aleatorios, que luego 
de su análisis por el método estadístico se propagarán al valor de la
constante $g$.

    \section{Resolución}
\subsection{Modelo entidad-relación}

El modelo entidad-relación que planteamos es el de la fig. \ref{fig:der}.
Asumimos las siguientes hipótesis:

\begin{itemize}
    \item La especialidad del Profesional es única. En caso de no serlo, se constata la
          más pertinente al trabajo que realiza.
    \item Pueden existir Áreas que todavía no tengan ni empleados ni director asignados.
    \item No todos los No Profesionales tienen instauradas franjas horarias desde que son
          ingresados a la base de datos.
    \item Los No Profesionales ocupan una única franja horaria. En caso de no existir,
          ésta se crea para asignársela.
    \item Los tipos de empleado son disjuntos. Un empleado no podrá ser de dos tipos a la
          vez.
    \item El tipo de Profesional es disjunto. Un empleado de Planta Permanente no podrá
          ser contratado y viceversa.
    \item El resultado de una auditoría es un booleano que puede o no venir acompañado de
          una descripción breve.
\end{itemize}

\begin{landscape}
    \begin{figure}[h]
        \centering
        \includegraphics[height=0.95\textheight]{/home/naza/Documents/drawing3.pdf}
        \caption{Diagrama Entidad-Relación}
        \label{fig:der}
    \end{figure}
\end{landscape}

Las restricciones sobre el esquema de la fig. \ref{fig:der} son las siguientes:

\begin{itemize}
    \item Todo empleado debe ser jerárquico, profesional o no profesional. A su vez entre
          jerárquico, profesional y no profesional se debe obtener a todos los empleados.
    \item Todo profesional debe ser contratado o de planta permanente. A su vez entre
          planta permanente y contratado se debe obtener a todos los empleados.
    \item Los empleados que acceden a áreas restringidas deberán estar capacitados para
          un nivel mayor o igual al del área.
    \item Dirección sólo puede tener empleados jerárquicos.
    \item Desarrolla sólo puede tener empleados contratados.
    \item Actúa\_en solo puede tener empleados de planta\_permanente.
    \item Toda área restringida es también un área.
    \item El nivel de seguridad de un área restringida no puede ser bajo.
    \item Sólo los profesionales tienen especialidades.
    \item Las direcciones de una misma área deben tener períodos que no se solapen.
    \item Las tareas que desarrolla un contratado no pueden solaparse en período.
    \item Un no\_profesional no puede ocupar dos franjas horarias que se solapen en
          periodo.
    \item Toda fecha se aloja como un \textit{timestamp}.
    \item En ningún atributo es posible asignar una fecha posterior a la fecha de carga
          de datos.
    \item No es posible que se produzca superposición entre dos períodos temporales del
          mismo tipo (fechas de inicio y finalización).
\end{itemize}

\section{Modelo lógico-relacional}

\newcommand{\pk}[1]{\underline{#1}}
\newcommand{\fk}[1]{\dashuline{#1}}
\newcommand{\ent}[1]{\uppercase{\texttt{#1}}}
\newcommand{\att}[1]{\texttt{\makefirstuc{#1}}}
\newcommand{\CK}[1]{CK=\{(#1)\}}
\newcommand{\FK}[1]{FK=\{(#1)\}}
\newcommand{\PK}[1]{PK=\{(#1)\}}
\newcommand{\fmin}[2]{$F_{\text{mín}}$: \{#1\} $\rightarrow$ \{#2\}}

\begin{itemize}
    \item \ent{area}(\pk{\att{nro\_area}}, \att{nombre},
          \att{es\_restringida})
          \begin{itemize}
              \item \ent{area}.\att{nro\_area} puede estar (o no) en
                    \ent{direccion}.\att{nro\_area}.
              \item \ent{area}.\att{nro\_area} puede estar (o no) en
                    \ent{tarea}.\att{nro\_area}.
              \item \ent{area}.\att{nro\_area} puede estar (o no) en
                    \ent{trabaja\_en}.\att{nro\_area}.
              \item \ent{area}.\att{nro\_area} puede estar (o no) en
                    \ent{planta\_permanente}.\att{nro\_area}.
              \item \ent{area}.\att{es\_restringida} es \{\texttt{true},
                    \texttt{false}\}.
              \item \CK{\att{nro\_area}}.
              \item \PK{\att{nro\_area}}.
              \item \FK{$\emptyset$}.
              \item Normalización:
                    \begin{itemize}
                        \item \fmin{\att{nro\_area}}{\att{nombre},
                                  \att{es\_restringida}}
                        \item \att{nro\_area} es superclave y esta a la
                              izquierda, por lo que está en BCNF.
                    \end{itemize}
          \end{itemize}

    \item \ent{area\_restringida}(\pk{\fk{\att{nro\_area}}}, \fk{\att{id\_seg}})
          \begin{itemize}
              \item \ent{area\_restringida}.\att{nro\_area} debe estar en
                    \ent{area}.\att{nro\_area}.
              \item \ent{area\_restringida}.\att{id\_seg} debe estar en
                    \ent{seguridad}.\att{id\_seg}.
              \item Para \ent{empleado}.\att{legajo} == \ent{area}.\att{legajo} y
                    \ent{area}.\att{nro\_area} en \- \ent{area\_restringida}.\att{nro\_area},
                    \ent{empleado}.\att{id\_seg} $\geq$ \ent{area\_restringida}.\att{id\_seg}.
              \item \CK{\att{nro\_area}}.
              \item \PK{\att{nro\_area}}.
              \item \FK{\att{nro\_area}, \att{id\_seg}}.
              \item Normalización:
                    \begin{itemize}
                        \item \fmin{\att{nro\_area}}{\att{seguridad}}.
                        \item BCNF(\{\att{nro\_area}\} es superclave de \ent{area\_restringida} y está del
                              lado izquierdo de su dependencia funcional.
                    \end{itemize}
          \end{itemize}

    \item \ent{auditoria}(\fk{\att{legajo}}, \fk{\att{id\_tarea}},
          \pk{\att{nro\_audit}}, \att{desc}, \att{fecha}, \att{aprobada})
          \begin{itemize}
              \item \ent{Auditoria}.\att{id\_tarea} debe estar
                    \ent{Desarrolla}.\att{id\_tarea}.
              \item \ent{Auditoria}.\att{legajo} debe estar en
                    \ent{Desarrolla}.\att{legajo}.
              \item \ent{Desarrolla}.\att{(legajo, id\_tarea)} puede estar (o
                    no) en \ent{Auditoria}.\att{(legajo, id\_tarea)}
              \item \ent{Auditoria}.\att{aprobada} es \{\texttt{true},
                    \texttt{false}\}.
              \item \CK{\att{legajo}, \att{id\_tarea}, \att{nro\_audit}}.
              \item \PK{\att{legajo}, \att{id\_tarea}, \att{nro\_audit}}.
              \item \FK{\att{legajo}, \att{id\_tarea}}.
              \item Normalización:
                    \begin{itemize}
                        \item \fmin{\att{legajo}, \att{id\_tarea},
                                  \att{nro\_audit}}{\att{desc}, \att{fecha},
                                  \att{aprobada}}.
                        \item BCNF(\{\att{legajo}, \att{id\_tarea}, \att{nro\_audit}\} es superclave de
                              \ent{auditoria} y está del lado izquierdo de su dependencia funcional).
                    \end{itemize}
          \end{itemize}

    \item \ent{contratado}(\pk{\fk{\att{legajo}}})
          \begin{itemize}
              \item \ent{contratado}.\att{legajo} debe estar en
                    \ent{profesional}.\att{legajo}.
              \item Para \ent{contratado}.\att{legajo} == \ent{profesional}.\att{legajo},
                    \ent{profesional}.\att{tipo\_prof} = 1.
              \item \CK{\att{legajo}}.
              \item \PK{\att{legajo}}.
              \item \FK{\att{legajo}}.
              \item Normalización:
                    \begin{itemize}
                        \item BCNF(no tiene dependencias funcionales no triviales, por lo que toda
                              dependencia cumple que la parte izquierda de la misma es superclave).
                    \end{itemize}
          \end{itemize}

    \item \ent{desarrolla}(\pk{\fk{\att{legajo}}},
          \pk{\fk{\att{id\_tarea}}}, \att{f\_ing}, \att{f\_fin})
          \begin{itemize}
              \item \ent{desarrolla}.\att{legajo} debe estar en
                    \ent{contratado}.\att{legajo}.
              \item \ent{desarrolla}.\att{id\_tarea} debe estar en
                    \ent{tarea}.\att{id\_tarea}.
              \item \ent{desarrolla}.\att{f\_fin} $>$
                    \ent{desarrolla}.\att{f\_ing}.
              \item \CK{\att{legajo}, \att{id\_tarea}}
              \item \PK{\att{legajo}, \att{id\_tarea}}
              \item \FK{\att{legajo}), (\att{id\_tarea}}
              \item Normalización:
                    \begin{itemize}
                        \item \fmin{\att{legajo}, \att{id\_tarea}}{\att{f\_ing},
                                  \att{f\_fin}}.
                        \item BCNF(\{\att{legajo}, \att{id\_tarea}\} es superclave de \ent{desarrolla} y está
                              del lado izquierdo de su dependencia funcional).
                    \end{itemize}
          \end{itemize}

    \item \ent{direccion}(\pk{\att{id\_dir}}, \att{f\_ini}, \att{f\_fin},
          \fk{\att{nro\_area}})
          \begin{itemize}
              % \item \ent{direccion}.\att{id\_dir} puede no estar en 
              %    \ent{ocupa}.\att{id\_dir}.
              \item \ent{direccion}.\att{nro\_area} debe estar en
                    \ent{area}.\att{nro\_area}.
              \item \ent{direccion}.\att{f\_fin} puede ser nulo si el director
                    sigue en su puesto.
              \item \CK{\att{id\_dir}), (\att{f\_ini}, \att{f\_fin},
                        \att{id\_area}}.
              \item \PK{\att{id\_dir}}.
              \item \FK{\att{nro\_area}}.
              \item Normalización:
                    \begin{itemize}
                        \item \fmin{(\att{id\_dir}),
                                  (\att{f\_ini}, \att{f\_fin}, \att{nro\_area})}
                              {(\att{f\_ini}, \att{f\_fin}, \att{nro\_area}),
                                  (\att{id\_dir})}.
                        \item Ambos elementos de la izquierda son superclave, por lo que están en BCNF.
                    \end{itemize}
          \end{itemize}

    \item \ent{durante}(\pk{\fk{\att{legajo}}}, \pk{\fk{\att{nro\_area}}}, \pk{\fk{\att{id\_franja}}})
          \begin{itemize}
              \item (\ent{durante}.\att{legajo},
                    \ent{durante}.\att{nro\_area}) puede no estar en
                    (\ent{trabaja\_en}.\att{legajo},
                    \ent{trabaja\_en}.\att{nro\_area}).
              \item \ent{durante}.\att{id\_franja} debe estar en
                    \ent{franja\_horaria}.\att{id\_franja}.
              \item \CK{\att{legajo}, \att{nro\_area}, \att{id\_franja}}.
              \item \PK{\att{legajo}, \att{nro\_area}, \att{id\_franja}}.
              \item \FK{\att{legajo}, \att{nro\_area}, \att{id\_franja}}.
              \item Normalización:
                    \begin{itemize}
                        \item BCNF(no tiene dependencias funcionales no triviales, por lo que toda
                              dependencia cumple que la parte izquierda de la misma es superclave)
                    \end{itemize}
          \end{itemize}

    \item \ent{Empleado}(\pk{\att{Legajo}}, \att{Nombre}, \att{Apellido},
          \att{F\_nac}, \att{Nro\_tel}, \att{Huella}, \att{Pwd}, \att{Tipo\_empl},
          \fk{\att{Id\_seg}})
          \begin{itemize}
              \item \ent{Empleado}.\att{Id\_seg} debe estar en
                    \ent{Seguridad}.\att{Id\_seg}.
              \item \ent{Empleado}.\att{Id\_seg} no puede ser nulo.
              \item \ent{Empleado}.\att{Tipo\_empl} es \{0 (\ent{profesional}), 1
                    (\ent{no\_profesional}), 2 (\ent{jerarqui\-co})\}.
              \item \ent{Empleado}.\att{Tipo\_empl} no puede ser nulo.
              \item \CK{\att{legajo}}.
              \item \PK{\att{legajo}}.
              \item \FK{\att{Id\_seg}}.
              \item Normalización:
                    \begin{itemize}
                        \item \fmin{\att{legajo}}{\att{Nombre}, \att{Apellido},
                                  \att{F\_nac}, \att{Nro\_tel}, \att{Huella},
                                  \att{Pwd}, \att{Tipo\_empl}, \att{Id\_seg}}.
                        \item BCNF({\att{legajo}} es superclave de \ent{empleado}, está del lado izquierdo de
                              su dependencia funcional).
                    \end{itemize}
          \end{itemize}

    \item \ent{entrada\_salida}(\pk{\fk{\att{legajo}}},
          \pk{\fk{\att{nro\_area}}}, \att{fecha}, \att{es\_entrada},
          \att{autorizado})
          \begin{itemize}
              \item \ent{entrada\_salida}.\att{legajo} debe estar en
                    \ent{empleado}.\att{legajo}.
              \item \ent{entrada\_salida}.\att{nro\_area} debe estar en
                    \ent{area\_restringida}.\att{nro\_area}.
              \item Si \ent{entrada\_salida}.\att{legajo} está en
                    \ent{planta\_permanente}.\att{legajo} entonces \att{legajo} debe determinar
                    funcionalmente a \att{nro\_area}.
              \item \ent{entrada\_salida}.\att{es\_entrada} es \{\texttt{true},
                    \texttt{false}\}.
              \item \ent{entrada\_salida}.\att{autorizado} es \{\texttt{true},
                    \texttt{false}\}.
              \item \CK{\att{legajo}, \att{nro\_area}}.
              \item \PK{\att{legajo}, \att{nro\_area}}.
              \item \FK{\att{legajo}), (\att{nro\_area}}.
              \item Normalización:
                    \begin{itemize}
                        \item \fmin{\att{legajo}, \att{nro\_area}}
                              {\att{fecha}, \att{es\_entrada}, \att{autorizado}}.
                        \item Como \{\att{legajo}, \att{nro\_area}\} es superclave y esta a la izquierda es
                              BCNF.
                    \end{itemize}
          \end{itemize}

    \item \ent{especialidad}(\pk{\att{id\_esp}}, \att{desc})
          \begin{itemize}
              \item \ent{especialidad}.\att{id\_esp} debe estar en
                    \ent{profesional}.\att{id\_esp}.
              \item \CK{\att{id\_esp}}.
              \item \PK{\att{id\_esp}}.
              \item \FK{$\emptyset$}.
              \item Normalización:
                    \begin{itemize}
                        \item \fmin{\att{id\_esp}}{\att{desc}}.
                        \item BCNF(\{\att{id\_esp}\} es superclave de \ent{especialidad} y está del lado
                              izquierdo de su dependencia funcional).
                    \end{itemize}
          \end{itemize}

    \item \ent{evento}(\fk{\att{nro\_area}}, \pk{\att{nro\_evento}},
          \att{desc}, \att{fecha})
          \begin{itemize}
              \item \ent{evento}.\att{nro\_area} debe estar en
                    \ent{area}.\att{nro\_area}.
              \item \ent{area}.\att{nro\_area} puede estar (o no) en
                    \ent{evento}.\att{nro\_area}.
              \item \CK{\att{nro\_evento}}.
              \item \PK{\att{nro\_evento}}.
              \item \FK{\att{nro\_area}}.
              \item Normalización:
                    \begin{itemize}
                        \item \fmin{\att{nro\_evento}}
                              {\att{nro\_area}, \att{desc}, \att{fecha}}.
                        \item BCNF(\att{nro\_evento} es superclave de \ent{evento} y está del lado izquierdo
                              de su dependencia funcional).
                    \end{itemize}
          \end{itemize}

    \item \ent{franja\_horaria}(\pk{\att{id\_franja}}, \att{hora\_entrada},
          \att{hora\_salida})
          \begin{itemize}
              \item \ent{franja\_horaria}.\att{id\_franja} debe estar en
                    \ent{trabaja\_en}.\att{id\_franja}.
              \item \ent{franja\_horaria}.\att{hora\_entrada} <
                    \ent{franja\_horaria}.\att{hora\_salida}.
              \item \CK{\att{id\_franja}\}, (\att{hora\_entrada},
                        \att{hora\_salida}}.
              \item \PK{\att{id\_franja}}
              \item \FK{$\emptyset$}
              \item Normalización:
                    \begin{itemize}
                        \item \fmin{\att{hora\_entrada},
                                  \att{hora\_salida}}{\att{id\_franja}}.
                        \item \fmin{\att{id\_franja}}{\att{hora\_entrada},
                                  \att{hora\_salida}}.
                        \item BCNF(\{\att{hora\_entrada}, \att{hora\_salida}\} e \att{id\_franja} son
                              superclave de \- \ent{franja\_horaria} y ambas están del lado izquierdo de su
                              dependencia funcional.
                    \end{itemize}
          \end{itemize}

    \item \ent{jerarquico}(\pk{\fk{\att{legajo}}})
          \begin{itemize}
              \item \ent{jerarquico}.\att{legajo} debe estar en
                    \ent{empleado}.\att{legajo}.
              \item Para \ent{jerarquico}.\att{legajo} == \ent{empleado}.\att{legajo},
                    \ent{empleado}.\att{tipo\_empl} = 2.
              \item \ent{jerarquico}.\att{legajo} puede no estar en
                    \ent{ocupa}.\att{legajo}.
              \item \CK{\att{legajo}}.
              \item \PK{\att{legajo}}.
              \item \FK{$\emptyset$}.
              \item Normalización:
                    \begin{itemize}
                        \item BCNF(no tiene dependencias funcionales no triviales, por lo que toda
                              dependencia cumple que la parte izquierda de la misma es superclave)
                    \end{itemize}
          \end{itemize}

    \item \ent{no\_profesional}(\pk{\fk{\att{legajo}}})
          \begin{itemize}
              \item \ent{no\_profesional}.\att{legajo} debe estar en
                    \ent{empleado}.\att{legajo}.
              \item Para \ent{no\_profesional}.\att{legajo} == \ent{empleado}.\att{legajo},
                    \ent{empleado}.\att{tipo\_empl} = 1.
              \item \ent{no\_profesional}.\att{legajo} debe estar en
                    \ent{trabaja\_en}.\att{legajo}.
              \item \CK{\att{legajo}}.
              \item \PK{\att{legajo}}.
              \item \FK{\att{legajo}}.
              \item Normalización:
                    \begin{itemize}
                        \item BCNF(no tiene dependencias funcionales no triviales, por lo que toda
                              dependencia cumple que la parte izquierda de la misma es superclave).
                    \end{itemize}
          \end{itemize}

    \item \ent{ocupa}(\pk{\fk{\att{legajo}}}, \pk{\fk{\att{id\_dir}}})
          \begin{itemize}
              \item \ent{ocupa}.\att{legajo} debe estar en
                    \ent{jerarquico}.\att{legajo}.
              \item \ent{ocupa}.\att{id\_dir} debe estar en
                    \ent{direccion}.\att{id\_dir}.
              \item \CK{(\att{legajo}, \att{id\_dir})}.
              \item \PK{(\att{legajo}, \att{id\_dir})}.
              \item \FK{\att{legajo}, \att{id\_dir}}.
              \item Normalización:
                    \begin{itemize}
                        \item BCNF(no tiene dependencias funcionales no triviales, por lo que toda
                              dependencia cumple que la parte izquierda de la misma es superclave)
                    \end{itemize}

          \end{itemize}

    \item \ent{planta\_permanente}(\pk{\fk{\att{legajo}}},
          \fk{\att{nro\_area}})
          \begin{itemize}
              \item \ent{planta\_permanente}.\att{legajo} debe estar en
                    \ent{profesional}.\att{legajo}.
              \item \ent{planta\_permanente}.\att{nro\_area} debe estar en
                    \ent{area}.\att{nro\_area}.
              \item Para \ent{planta\_permanente}.\att{legajo} == \ent{profesional}.\att{legajo},
                    \ent{profesional}.\att{tipo\_prof} = 0.
              \item \ent{empleado}.\att{id\_seg} de
                    \ent{planta\_permanente}.\att{legajo} $\geq$
                    \ent{area\_restringida}.\att{seguridad} de
                    \ent{planta\_permanente}.\att{nro\_area}.
              \item \CK{\att{legajo}}.
              \item \PK{\att{legajo}}.
              \item \FK{\att{nro\_area}}.
              \item Normalización:
                    \begin{itemize}
                        \item \fmin{\att{legajo}}{\att{nro\_area}}.
                        \item BCNF(\{\att{legajo}\} es superclave de \ent{planta\_permanente} y está del lado
                              izquierdo de su dependencia funcional).
                    \end{itemize}
          \end{itemize}

    \item \ent{profesional}(\pk{\fk{\att{legajo}}}, \att{tipo\_prof},
          \fk{\att{id\_esp}})
          \begin{itemize}
              \item \ent{profesional}.\att{tipo\_prof} es \{0
                    (\ent{planta\_permanente}), 1 (\ent{contratado})\}.
              \item \ent{profesional}.\att{legajo} debe estar en
                    \ent{empleado}.\att{legajo}.
              \item \ent{profesional}.\att{id\_esp} debe estar en
                    \ent{especialidad}.\att{id\_esp}.
              \item Para \ent{profesional}.\att{legajo} == \ent{empleado}.\att{legajo},
                    \ent{empleado}.\att{tipo\_emp} = 0.
              \item \CK{\att{legajo}}.
              \item \PK{\att{legajo}}.
              \item \FK{\att{id\_seg}}.
              \item Normalización:
                    \begin{itemize}
                        \item \fmin{\att{legajo}}{\att{tipo\_prof},
                                  \att{id\_esp}}.
                        \item BCNF(\{\att{legajo}\} es superclave de \ent{profesional} y está del lado
                              izquierdo de su dependencia funcional.
                    \end{itemize}
          \end{itemize}

    \item \ent{seguridad}(\pk{\att{Id\_seg}}, \att{Desc}, \att{Nivel})
          \begin{itemize}
              \item \ent{seguridad}.\att{id\_seg} debe estar en
                    \ent{area\_restringida}.\att{id\_seg}.
              \item Los niveles de seguridad se ordenan según \att{nivel}. Con esto se permite
                    añadir nuevos niveles, incluso intermedios, sin modificar el resto de las
                    instancias.
              \item \CK{\att{id\_seg}}.
              \item \PK{\att{id\_seg}}.
              \item \FK{$\emptyset$}.
              \item Normalización:
                    \begin{itemize}
                        \item \fmin{\att{id\_seg}}{\att{desc}, \att{nivel}}.
                        \item BCNF(\{\att{id\_seg}\} es superclave de \ent{seguridad} y está del lado
                              izquierdo de su dependencia funcional.
                    \end{itemize}
          \end{itemize}

    \item \ent{tarea}(\pk{\att{id\_tarea}}, \att{desc},
          \fk{\att{nro\_area}})
          \begin{itemize}
              \item \ent{tarea}.\att{nro\_area} debe estar en
                    \ent{area}.\att{nro\_area}.
              \item \ent{area}.\att{nro\_area} puede estar (o no) en
                    \ent{tarea}.\att{nro\_area} (esto va en \ent{area}).
              \item \CK{\att{id\_tarea}}.
              \item \PK{\att{id\_tarea}}.
              \item \FK{\att{nro\_tarea}}.
              \item Normalización:
                    \begin{itemize}
                        \item \fmin{\att{id\_tarea}}{\att{desc}, \att{nro\_area}}.
                        \item BCNF(\{\att{id\_tarea}\} es superclave de \ent{tarea} y está del lado izquierdo
                              de su dependencia funcional.
                    \end{itemize}
          \end{itemize}

    \item \ent{trabaja\_en}(\pk{\fk{\att{legajo}}}, \pk{\fk{\att{nro\_area}}},
          \fk{\att{id\_franja}})
          \begin{itemize}
              \item \ent{Trabaja\_en}.(\att{legajo}, \att{nro\_area}) debe
                    estar en \ent{No\_profesional}.\att{legajo}
              \item \ent{Trabaja\_en}.\att{nro\_area} debe estar en
                    \ent{Area}.\att{nro\_area}.
              \item \ent{Trabaja\_en}.\att{id\_franja} debe estar en
                    \ent{Franja\_horaria}.\att{id\_franja} o ser nulo.
              \item \ent{Empleado}.\att{id\_seg} de
                    \ent{trabaja\_en}.\att{legajo} $\geq$
                    \ent{area}.\att{seguridad} de
                    \ent{trabaja\_en}.\att{nro\_area}
              \item Si \ent{Trabaja\_en}.\att{legajo} $=$ \ent{Trabaja\_en}.\att{legajo} con
                    \ent{Trabaja\_en} $\neq$ \ent{Trabaja\_en}, $\Longrightarrow$
                    \ent{Trabaja\_en}.\att{franja\_horaria} $\neq$
                    \ent{Trabaja\_en}.\att{franja\_horaria}.
              \item \CK{\att{legajo}, \att{nro\_area}}.
              \item \PK{\att{legajo}, \att{nro\_area}}.
              \item \FK{\att{legajo}), (\att{nro\_area}), (\att{id\_franja}}.
              \item Normalización:
                    \begin{itemize}
                        \item \fmin{(\att{legajo}, \att{id\_tarea}),
                                  (\att{legajo}, \att{id\_franja})}
                              {(\att{id\_franja}), (\att{nro\_area})}
                        \item BCNF(\{\att{legajo}, \att{nro\_area}\} y \{\att{legajo}, \att{id\_franja}\} son
                              superclave de \ent{trabaja\_en} y ambas están del lado izquierdo de su
                              dependencia funcional)
                    \end{itemize}
          \end{itemize}

\end{itemize}
    
    \newpage

\end{document}
